\begin{table}[H]
\centering
\scriptsize
\caption{Findings register of the credit risk study. Severities are PROPOSED and must be confirmed by the validator.}
\label{tab:C6_findings}
\setlength{\tabcolsep}{3pt}
\begin{tabular}{>{\raggedright\arraybackslash}p{0.75cm}>{\raggedright\arraybackslash}p{6.6cm}>{\raggedright\arraybackslash}p{1.85cm}>{\raggedright\arraybackslash}p{2.4cm}>{\raggedright\arraybackslash}p{3.5cm}}
\toprule
ID & Finding & Severity & Evidence & Proposed remediation \\
\midrule
C-F1 & Calibration tests that assume independent defaults reject correct PDs in a single period: the Hosmer--Lemeshow test (chi-square with $g = 10$ degrees of freedom, as the true PD is not fitted on the sample) rejects the TRUE PD in 82.5 $\pm$ 2.7\,\% of the periods with asset correlation 0.12, against 5.5 $\pm$ 1.6\,\% (its nominal 5\,\%) with independent defaults; a one-year default rate carries the systematic factor. The per-grade binomial and Jeffreys tests are a multiple test over about 8 grades: they reject already with independent defaults (27.5 $\pm$ 3.2 / 34.5 $\pm$ 3.4\,\%), and correlation adds little (36.0 $\pm$ 3.4 / 38.5 $\pm$ 3.4\,\%). & PROPOSED: High & Tab.~\ref{tab:C2_calibration}; Fig.~\ref{fig:fig20_calibration_drift} & Interpret one-period calibration tests with correlation-adjusted critical values or over a cycle; control the family-wise level of per-grade tests (e.g. Bonferroni / Holm). \\
C-F2 & The MLE standard errors ignore the default correlation within a period: over the 200 replications with 50000 borrower-years (5 periods) the macro coefficient (true -0.45) has a spread of 1.026 (sd), 12 times its mean MLE standard error (0.087): the macro effect is identified from only 5 period-level observations, each shifted by the unobserved systematic factor. & PROPOSED: Medium & Tab.~\ref{tab:C1_pd_model} & Cluster standard errors by period, or estimate macro sensitivities on long histories. \\
C-F3 & In a downturn (macro at -2 sd) a model without the macro variable predicts 1.65 $\pm$ 0.05\,\% against a realised 3.65 $\pm$ 0.20\,\%. The model with the macro variable, estimated from 8 periods, overshoots on average (5.46 $\pm$ 0.67\,\%; true PD 3.55\,\%), but its median over the replications is 3.22\,\%: the overshoot of the mean is a Jensen effect - the noisy macro coefficient enters through the exponential, so the mean prediction is biased upwards; the median fit is close to the true PD. & PROPOSED: High & Tab.~\ref{tab:C2_calibration}; Fig.~\ref{fig:fig20_calibration_drift} & Monitor the calibration through the cycle; constrain macro sensitivities (expert bounds, longer data). \\
C-F4 & Small samples cost discrimination and stability: with 500 borrower-years 74 of 200 fits fail with separation, and the successful fits lose 0.072 AUC against the true PD; with 50000 the gap is 0.018. Failed fits are excluded, so the small-sample averages are conditional on success and flatter the model. & PROPOSED: Medium & Tab.~\ref{tab:C1_pd_model}; Fig.~\ref{fig:fig19_pd_sample_size} & Pool data or simplify the model for low-default portfolios; never hide failed fits. \\
C-F5 & The ECL depends strongly on two judgements: the coverage ranges from 0.97\,\% (no downside weight, SICR $\times$5) to 1.84\,\% (downside weight 60\,\%, SICR $\times$1.5). The SICR threshold moves the stage-2 share from 2 to 30\,\% and dominates; the downside weight changes the measurement only (staging compares PDs on the same basis, without the scenario overlay). The stage-2 transfer makes a single loan's provision jump by a multiple. & PROPOSED: Medium & Tab.~\ref{tab:C3_ecl}; Fig.~\ref{fig:fig21_ecl_cliff} & Disclose the sensitivity to the scenario weights and the SICR threshold; back-test the staging. \\
C-F6 & The IRB capital matches the one-factor 99.9\,\% loss of a granular portfolio (ratio 0.993) but understates a concentrated one: with 30 names the simulated unexpected loss is 1.69 times the sum of K. & PROPOSED: Medium & Tab.~\ref{tab:C4_irb}; Fig.~\ref{fig:fig22_irb_vs_mc} & Add a concentration charge (Pillar 2 / granularity adjustment) for portfolios with few names. \\
C-F7 & The IRB formula reproduces the illustrative risk weights of Basel II Annex 5 / CRE99 exactly, and the PD floor (CRE32.4), the M bounds (CRE32.46) and the CRR-only 1.06 scaling are confirmed; still used as recalled: the IFRS 9 paragraph numbers other than B5.5.13 and the ECB form of the Jeffreys test. & PROPOSED: Low & Tab.~\ref{tab:C5_assumptions} & Verify the remaining IFRS 9 references and the ECB instructions before any use beyond this showcase. \\
\bottomrule
\end{tabular}
\end{table}
